% Generated draft: requires mathematical and source review.
\documentclass[11pt]{article}
\usepackage[margin=25mm]{geometry}
\usepackage{fontspec}
\setmainfont{Latin Modern Roman}
\usepackage{amsmath,amssymb,mathtools,unicode-math}
\setmathfont{Latin Modern Math}
\usepackage{enumitem,needspace,xurl,hyperref}
\setlength{\parindent}{0pt}
\setlength{\parskip}{5pt}
\setlength{\emergencystretch}{3em}
\newcommand{\N}{\mathbb{N}}
\newcommand{\Z}{\mathbb{Z}}
\newcommand{\Q}{\mathbb{Q}}
\newcommand{\R}{\mathbb{R}}
\newcommand{\C}{\mathbb{C}}
\newcommand{\F}{\mathbb{F}}
\newcommand{\A}{\mathbb{A}}
\newcommand{\PP}{\mathbb{P}}
\newcommand{\E}{\mathbb{E}}
\newcommand{\eps}{\varepsilon}
\newcommand{\dd}{\,\mathrm{d}}
\begin{document}
\section*{UnsolvedMath ID: 1200004}
\textbf{Source problem code:} catalogue entry
\subsection{Definitions and mathematical statement}
Let $p$ be a prime, let $\Q_p$ be the field of $p$-adic numbers, and set
\[
 G=\operatorname{SL}_2(\Q_p)=\{A\in M_2(\Q_p):\det A=1\}.
\]
A possible intended matrix convention calls $A\in G$ parabolic when $\operatorname{tr}(A)\in\{2,-2\}$ and $A\notin\{I,-I\}$, equivalently when it is a noncentral matrix with a repeated eigenvalue. The record does not define ``parabolic''; this convention must be checked, particularly because terminology for actions on the Bruhat--Tits tree can differ.

Let $\Gamma\leq G$ be countable and contain no parabolic elements. For $g\in G$, $\langle\Gamma,g\rangle$ denotes the smallest subgroup containing $\Gamma$ and $g$. The source asks whether adjoining a random $g$ almost surely preserves the absence of parabolics. Because $G$ is noncompact, Haar measure on $G$ is not itself a probability distribution. The accompanying research summary makes the randomness precise as follows: let $m_G$ be Haar measure and let $\mu$ be \emph{any} Borel probability measure satisfying $\mu\ll m_G$. Show that
\[
 \mu\bigl(\{g\in G:\langle\Gamma,g\rangle
                  \text{ contains no parabolic element}\}\bigr)=1.
\]
Under the matrix convention above, the event can be written explicitly as
\[
 \forall A\in\langle\Gamma,g\rangle,
 \quad\operatorname{tr}(A)\in\{2,-2\}\ \Longrightarrow\ A\in\{I,-I\}.
\]
No discreteness assumption on $\Gamma$ is supplied.

The registry marks this solved by its own contributor. Its summary claims a proof for every such absolutely continuous probability measure. It describes a purported new algebraic lemma about generalized word maps on $\operatorname{SL}_2$ over algebraically closed fields of characteristic zero: identically $\pm2$ trace is said to force formal collapse to the constant $\pm I$. The summary says the no-parabolics assumption forces an endpoint product $a_0\cdots a_k=\pm I$, and that induction on the number of variable factors uses the top two coefficients of $\operatorname{tr}(w(I+tN))$ over the nilpotent cone. No proof or precise definition of formal collapse is supplied. These are attributed proof claims, with novelty qualified by the contributor's own verification; they are not certified results in this draft.
\subsection{Short English statement}
Given a countable subgroup of $\operatorname{SL}_2(\Q_p)$ with no parabolic elements, must adjoining an independently chosen element with any Haar-absolutely-continuous probability law almost surely preserve that property? The summary specifies this probability convention and claims an unpublished proof; the meaning of parabolic and that proof require review.
\subsection{Sources}
\begin{itemize}
\item Local UnsolvedMath registry, permanent ID 1200004, source code catalogue entry, ``Some Questions --- Question 4''; supplied status solved, classification SOLVED-BY-YOU, updated 2026-08-06. The proof and novelty claims are those of the supplied research summary.
\item Abert, \emph{Some questions}, author list dated 2010, Question 4, \url{https://www.renyi.hu/~abert/questions.pdf}. The record reports manual transcription and access on 2026-07-29; it says that the 2010 list was inspected and current status was not independently verified. Status and rights: NEEDS\_REVIEW; public author-hosted PDF.
\item No separate publication, proof manuscript, or external citation for the claimed algebraic lemma or claimed solution is supplied.
\end{itemize}
\end{document}
